\section*{Capítulo \ref{ch:g12:diversification-of-life} --- A diversificação dos seres vivos}

\begin{solution}{exo:g12:diversification-of-life:1}
Duplicação e divergência de \omterm{def:g10:universal-dna:gene}{genes}; \omterm{prop:g11:antibiotic-resistance:variation}{transferência horizontal de genes};
hibridação com poliploidia; mudanças na regulação dos \omterm{prop:g12:diversification-of-life:development}{genes do
desenvolvimento}; \omterm{def:g12:diversification-of-life:symbiosis}{simbiose}; comportamento transmitido.
\end{solution}

\begin{solution}{exo:g12:diversification-of-life:2}
Um conjunto de \omterm{def:g10:universal-dna:gene}{genes} aparentados descendentes de um \omterm{def:g10:universal-dna:gene}{gene} ancestral por
duplicações sucessivas, com cada cópia divergindo depois por \omterm{def:g11:mutations:mutation}{mutação}
para uma função distinta (\omterm{def:g11:the-eye:photoreceptors}{opsinas}, globinas).
\end{solution}

\begin{solution}{exo:g12:diversification-of-life:3}
Os \omterm{def:g10:universal-dna:information}{cromossomos} dele, um conjunto de cada \omterm{def:g10:biodiversity-scales:species}{espécie} genitora, não têm
homólogos para se emparelhar na \omterm{def:g12:meiosis-genetic-shuffling:meiosis}{meiose}, e assim os gametas ficam
desequilibrados. A duplicação dá a cada \omterm{def:g11:cell-cycle-mitosis:chromosome}{cromossomo} um parceiro
idêntico: o emparelhamento e a \omterm{def:g12:meiosis-genetic-shuffling:meiosis}{meiose} ficam regulares e a planta é fértil.
\end{solution}

\begin{solution}{exo:g12:diversification-of-life:4}
\omterm{def:g10:universal-dna:information}{DNA} circular próprio; \omterm{def:g10:cells-common-unit:organelle}{ribossomos} do tipo bacteriano sensíveis aos
\omterm{def:g11:antibiotic-resistance:antibiotic}{antibióticos} antibacterianos; multiplicação por divisão; uma membrana
dupla; sequências mais próximas de um grupo de bactérias vivas.
\end{solution}

\begin{solution}{exo:g12:diversification-of-life:5}
Um comportamento aprendido com os outros membros do grupo e passado
adiante sem \omterm{def:g10:universal-dna:gene}{genes}: a quebra de nozes com pedras em algumas comunidades
de chimpanzés e em outras não.
\end{solution}

\begin{solution}{exo:g12:diversification-of-life:6}
$14 + 7 = 21$ \omterm{def:g10:universal-dna:information}{cromossomos}: 7 A do einkorn, 7 A e 7 B do emmer. Sete
pares (A com A) conseguem se formar; os 7 \omterm{def:g10:universal-dna:information}{cromossomos} B não têm
parceiro e se distribuem ao acaso, e assim os gametas ficam
desequilibrados: estéril.
\end{solution}

\begin{solution}{exo:g12:diversification-of-life:7}
A separação $\alpha$--$\beta$ (50\% idênticos) é mais antiga que a
separação $\beta$--$\gamma$ (80\%): $\alpha$ se ramifica primeiro, e
depois $\beta$ e $\gamma$ se separam um do outro.
\end{solution}

\begin{solution}{exo:g12:diversification-of-life:8}
Os \omterm{def:g10:universal-dna:gene}{genes} dos membros estão presentes, mas o sinal que os liga no flanco
do embrião não é dado, ou é dado e depois interrompido; a diferença
está na regulação da expressão, não nos \omterm{def:g10:universal-dna:gene}{genes}.
\end{solution}

\begin{solution}{exo:g12:diversification-of-life:9}
Na quantidade e no momento de um sinal de crescimento no bico
embrionário: uma diferença de regulação. Uma mudança de regulação
precisa apenas de uma pequena \omterm{def:g11:mutations:mutation}{mutação} numa sequência de controle,
enquanto uma \omterm{def:g11:gene-expression:protein}{proteína} nova precisaria de muitas; ela pode surgir e ser
selecionada em alguns milhares de gerações.
\end{solution}

\begin{solution}{exo:g12:diversification-of-life:10}
A forma dele, a capacidade dele de colonizar rocha nua, a resistência
dele à seca e o crescimento lento dele não pertencem a nenhum dos
parceiros sozinho; eles emergem da associação. O líquen é um organismo
composto, com propriedades próprias.
\end{solution}

\begin{solution}{exo:g12:diversification-of-life:11}
Comportamento transmitido: os materiais estão disponíveis dos dois
lados, as populações são próximas geneticamente, e a diferença
acompanha o rio --- uma barreira ao contato e, portanto, ao
aprendizado, não aos \omterm{def:g10:universal-dna:gene}{genes} nem aos recursos.
\end{solution}

\begin{solution}{exo:g12:diversification-of-life:12}
Um vírus inseriu o \omterm{def:g10:universal-dna:gene}{gene} dele no genoma da \omterm{def:g10:cells-common-unit:cell}{célula} germinativa de um
ancestral dos primatas há algumas dezenas de milhões de anos; o \omterm{def:g10:universal-dna:gene}{gene}
foi herdado e conservado. Teste: o mesmo \omterm{def:g10:universal-dna:gene}{gene} deve estar na mesma
posição cromossômica em todos os primatas e estar ausente dessa posição
nos outros mamíferos; a sequência dele deve ser a mais próxima da do vírus.
\end{solution}

\begin{solution}{exo:g12:diversification-of-life:13}
Os \omterm{def:g10:cells-common-unit:organelle}{ribossomos} mitocondriais são do tipo bacteriano e ligam os mesmos
\omterm{def:g11:antibiotic-resistance:antibiotic}{antibióticos}: em doses altas o medicamento retarda a síntese de
\omterm{def:g11:gene-expression:protein}{proteínas} das \omterm{def:g10:cells-common-unit:organelle}{mitocôndrias}. A sensibilidade é herdada do ancestral
bacteriano das \omterm{def:g10:cells-common-unit:organelle}{mitocôndrias}.
\end{solution}

\begin{solution}{exo:g12:diversification-of-life:14}
O número de \omterm{def:g10:universal-dna:information}{cromossomos} do \omterm{prop:g12:diversification-of-life:polyploidy}{poliploide} não bate mais com o de nenhum dos
genitores, e assim os cruzamentos com eles dão descendentes estéreis: o
isolamento é imediato. Um \omterm{def:g10:universal-dna:gene}{alelo} novo precisa se espalhar por uma
população ao longo de muitas gerações antes de distinguir um grupo. Uma
\omterm{def:g10:biodiversity-scales:species}{espécie} de planta pode, portanto, surgir numa única geração.
\end{solution}

\begin{solution}{exo:g12:diversification-of-life:15}
No fundo, toda sequência nova --- uma cópia duplicada que diverge, uma
mudança de regulação, os \omterm{def:g10:universal-dna:gene}{genes} de um plasmídeo transferido --- começou
como uma \omterm{def:g11:mutations:mutation}{mutação} em algum lugar. Mas a duplicação, a transferência e a
poliploidia movem e multiplicam \omterm{def:g10:universal-dna:gene}{genes} e genomas inteiros de uma vez; a
mudança de regulação produz formas novas sem \omterm{def:g11:gene-expression:protein}{proteínas} novas; a
\omterm{def:g12:diversification-of-life:symbiosis}{simbiose} combina genomas de \omterm{def:g10:biodiversity-scales:species}{espécies} diferentes; e a cultura transmite
variação sem \omterm{def:g10:universal-dna:gene}{gene} nenhum. A \omterm{def:g11:mutations:mutation}{mutação} fornece as letras; esses processos
reorganizam páginas e livros inteiros.
\end{solution}

\begin{solution}{pb:g12:diversification-of-life:1}
\textbf{1.} Emmer 28, trigo comum 42.

\textbf{2.} Nenhum: os \omterm{def:g10:universal-dna:information}{cromossomos} A e B não são homólogos. Os gametas
recebem sortimentos aleatórios dos 14, quase nunca um conjunto completo.

\textbf{3.} Com 28 \omterm{def:g10:universal-dna:information}{cromossomos}, cada A tem um parceiro A idêntico e
cada B um B: 14 pares, \omterm{def:g12:meiosis-genetic-shuffling:meiosis}{meiose} regular, gametas equilibrados de 14.

\textbf{4.} Nenhum consegue se emparelhar: 7 \omterm{def:g10:universal-dna:information}{cromossomos} A, 7 B e 7 D,
cada um sem homólogo. Estéril.

\textbf{5.} 21 pares; um grão de pólen carrega 21.

\textbf{6.} As outras duas cópias continuam fornecendo a \omterm{def:g11:gene-expression:protein}{proteína}. As
cópias redundantes ficam livres para acumular mutações e divergir ---
duplicação na escala de um genoma inteiro.

\textbf{7.} Nos campos da região onde o emmer era cultivado, entre
10\,000 e 8\,000 anos atrás, com o aegilops silvestre 2 crescendo como
erva daninha no meio da lavoura.

\textbf{8.} Que os dois passos --- hibridação e duplicação --- bastam
para produzir o trigo comum a partir dos genitores dele: a história
pode ser refeita.

\textbf{9.} Sim: ele não consegue produzir descendentes férteis com o
einkorn, e assim, pela definição, é uma \omterm{def:g10:biodiversity-scales:species}{espécie} distinta, embora
descenda dele.

\textbf{10.} Os \omterm{def:g10:universal-dna:information}{cromossomos} D do aegilops 2 são homólogos ao conjunto D
do trigo e se emparelham com ele, e assim um cruzamento seguido de
retrocruzamentos consegue introduzir o \omterm{def:g10:universal-dna:gene}{gene}; os \omterm{def:g10:universal-dna:information}{cromossomos} de uma
planta sem parentesco não se emparelhariam de jeito nenhum.

\textbf{11.} Os \omterm{def:g10:universal-dna:gene}{genes} em si quase não diferem; a diferença entre os
corpos deve estar no modo como os \omterm{def:g10:universal-dna:gene}{genes} são usados.

\textbf{12.} Uma mudança na regulação da expressão --- a extensão da
região ao longo do eixo em que o \omterm{def:g10:universal-dna:gene}{gene} fica ativo.

\textbf{13.} Os \omterm{def:g10:universal-dna:gene}{genes} dos membros estão presentes e começam a agir; o
sinal que sustenta o crescimento do membro está faltando, e assim os
brotos param. De novo a regulação, e não a perda dos \omterm{def:g10:universal-dna:gene}{genes}.

\textbf{14.} Ele mostra que mudar onde o \omterm{def:g10:universal-dna:gene}{gene} é expresso basta para
mudar a região do corpo: a correlação entre expressão e forma é causal.

\textbf{15.} A mudança do trigo comum acrescentou conjuntos inteiros ao
genoma (o tamanho dele); a da píton mudou onde os \omterm{def:g10:universal-dna:gene}{genes} compartilhados
são usados (a regulação deles).

\textbf{16.} \omterm{def:g10:universal-dna:information}{DNA} circular próprio (esperado: nenhum, todos os \omterm{def:g10:universal-dna:gene}{genes}
nucleares); \omterm{def:g10:cells-common-unit:organelle}{ribossomos} bacterianos (esperado: \omterm{def:g10:cells-common-unit:organelle}{ribossomos}
eucarióticos); multiplicação por divisão (esperado: montagem a partir
de peças); membrana dupla (esperado: uma membrana única como a das
outras \omterm{def:g10:cells-common-unit:organelle}{organelas}). E os \omterm{def:g10:universal-dna:gene}{genes} deles se parecem com \omterm{def:g10:universal-dna:gene}{genes} bacterianos
(esperado: \omterm{def:g10:universal-dna:gene}{genes} do tipo nuclear).

\textbf{17.} No \omterm{def:g10:cells-common-unit:organelle}{núcleo}: com o tempo, a maioria dos \omterm{def:g10:universal-dna:gene}{genes} da bactéria
foi transferida para o genoma nuclear, e o genoma próprio da \omterm{def:g10:cells-common-unit:organelle}{organela}
encolheu até virar um resquício.

\textbf{18.} Os \omterm{def:g10:universal-dna:gene}{genes} mitocondriais passam apenas de mãe para filho; o
\omterm{def:g10:universal-dna:information}{DNA} mitocondrial de uma pessoa traça a linha materna ininterrupta.

\textbf{19.} A \omterm{def:g10:cells-common-unit:organelle}{mitocôndria} cedeu a maior parte dos \omterm{def:g10:universal-dna:gene}{genes} dela ao
hospedeiro e não consegue existir fora da \omterm{def:g10:cells-common-unit:cell}{célula}; os parceiros do
líquen conservam os genomas deles e podem, com dificuldade, ser separados.

\textbf{20.} Einkorn (14, A), aegilops 1 (14, B), aegilops 2 (14, D):
42 \omterm{def:g10:universal-dna:information}{cromossomos}; ``regulação''; a \omterm{def:g10:cells-common-unit:organelle}{mitocôndria}.
\end{solution}
